%% kapa: exam=1 round=29 year=2018 period=1 form=A subject=civil_law subject_ko=민법 seq=24 number=24
%% source: Q-Net 감정평가사 기출문제 (cst00309) artlSeq=5202808 file=extracted/1008/5202808/2018년 제29회 감정평가사 1차시험 기출문제/2018년 제29회 감정평가사 1차시험 기출문제/1교시 A형.pdf
%% conversion: kordoc 4.24.51 -> tools/gen1.py (auto); review=proofread-sonnet; flags=none
\begin{question}{24}
유치권에 관한 설명으로 옳지 않은 것은? (다툼이 있으면 판례에 따름)
\begin{choices}
  \item 유치권의 행사는 피담보채권의 소멸시효의 진행에 영향을 미치지 아니한다.
  \item 유치권자는 피담보채권 전부의 변제를 받을 때까지 유치물 전부에 대하여 그 권리를 행사할 수 있다.
  \item 근저당권설정 후 그 실행에 따른 경매로 인한 압류의 효력이 발생하기 전에 취득한 유치권으로 경매절차의 매수인에게 대항할 수 없다.
  \item 피담보채권의 채무자를 직접점유자로 하여 채권자가 간접점유하는 경우에 유치권은 성립하지 않는다.
  \item 유치권자는 경매로 인한 매수인에 대하여 그 피담보채권의 변제가 있을 때까지 유치목적물의 인도를 거절할 수 있을 뿐, 그 피담보채권의 변제를 청구할 수는 없다.
\end{choices}
\end{question}
\answer{3}
